\newcommand{\Qed}{\begin{flushright}\qed\end{flushright}}
\newcommand{\parinn}{\setlength{\parindent}{1cm}}
\newcommand{\parinf}{\setlength{\parindent}{0cm}}
\newcommand{\del}[2]{\frac{\partial #1}{\partial #2}}
\newcommand{\Del}[3]{\frac{\partial^{#1} #2}{\partial^{#1} #3}}
\newcommand{\deld}[2]{\dfrac{\partial #1}{\partial #2}}
\newcommand{\Deld}[3]{\dfrac{\partial^{#1} #2}{\partial^{#1} #3}}
\newcommand{\uin}{\mathbin{\rotatebox[origin=c]{90}{$\in$}}}
\newcommand{\usubset}{\mathbin{\rotatebox[origin=c]{90}{$\subset$}}}
\newcommand{\lt}{\left}
\newcommand{\rt}{\right}
\newcommand{\exs}{\exists}
\newcommand{\st}{\strut}
\newcommand{\dps}[1]{\displaystyle{#1}}
\newcommand{\la}{\langle}
\newcommand{\ra}{\rangle}
\newcommand{\cls}[1]{\textsc{#1}}
\newcommand{\prb}[1]{\textsc{#1}}
\newcommand{\comb}[2]{\left(\begin{matrix}
		#1\\ #2
\end{matrix}\right)}
%\newcommand[2]{\quotient}{\faktor{#1}{#2}}
\newcommand\quotient[2]{
	\mathchoice
	{% \displaystyle
		\text{\raise1ex\hbox{$#1$}\Big/\lower1ex\hbox{$#2$}}%
	}
	{% \textstyle
		#1\,/\,#2
	}
	{% \scriptstyle
		#1\,/\,#2
	}
	{% \scriptscriptstyle  
		#1\,/\,#2
	}
}

\newcommand{\tensor}{\otimes}
\newcommand{\xor}{\oplus}

\newcommand{\sol}[1]{\begin{solution}#1\end{solution}}
\newcommand{\solve}[1]{\setlength{\parindent}{0cm}\textbf{\textit{Solution: }}\setlength{\parindent}{1cm}#1 \hfill $\blacksquare$}
\newcommand{\mat}[1]{\left[\begin{matrix}#1\end{matrix}\right]}
\newcommand{\matr}[1]{\begin{matrix}#1\end{matrix}}
\newcommand{\matp}[1]{\lt(\begin{matrix}#1\end{matrix}\rt)}
\newcommand{\detmat}[1]{\lt|\begin{matrix}#1\end{matrix}\rt|}
\newcommand\numberthis{\addtocounter{equation}{1}\tag{\theequation}}
\newcommand{\handout}[3]{
	\noindent
	\begin{center}
		\framebox{
			\vbox{
				\hbox to 6.5in { {\bf Complexity Theory I } \hfill Jan -- May, 2023 }
				\vspace{4mm}
				\hbox to 6.5in { {\Large \hfill #1  \hfill} }
				\vspace{2mm}
				\hbox to 6.5in { {\em #2 \hfill #3} }
			}
		}
	\end{center}
	\vspace*{4mm}
}

\newcommand{\lecture}[3]{\handout{Lecture #1}{Lecturer: #2}{Scribe:	#3}}

\let\marvosymLightning\Lightning
\newcommand{\ctr}{\text{\marvosymLightning}\xspace} % Requires marvosym package

\newcommand{\ov}[1]{\overline{#1}}
\newcommand{\by}[1]{&[\text{By #1}]}
\newcommand{\byt}[2][]{\by{\thmref[#1]{thm:#2}}}
\newcommand{\byl}[2][]{\by{\lemref[#1]{lem:#2}}}
\newcommand{\byc}[2][]{\by{\corref[#1]{cor:#2}}}
\newcommand{\byo}[2][]{\by{\obsref[#1]{obs:#2}}}
\newcommand{\bye}[1]{\by{\eqref{eq:#1}}}
\newcommand{\bm}[1]{\boldsymbol{#1}}
\newcommand{\Leg}[2]{\genfrac{(}{)}{}{0}{#1}{#2}}
\newcommand{\Xq}{\bm{\mathcal{X}}_q}
\newcommand{\Mq}{\bm{\mathcal{M}}_q}
\newcommand{\lmq}{\bm{\Lambda}_q}
\newcommand{\Xp}{\bm{\mathcal{X}}_p}
\newcommand{\Mp}{\bm{\mathcal{M}}_p}
\newcommand{\thmref}[2][]{%
	\hyperref[th:#2]{Theorem~\ref{th:#2}#1}%
}
\newcommand{\propref}[2][]{%
	\hyperref[th:#2]{Proposition~\ref{th:#2}#1}%
}
\newcommand{\lemref}[2][]{%
	\hyperref[th:#2]{Lemma~\ref{th:#2}#1}%
}
\newcommand{\corref}[2][]{%
	\hyperref[th:#2]{Corollary~\ref{th:#2}#1}%
}
\newcommand{\obsref}[2][]{%
  \hyperref[#2]{Observation~\ref{#2}#1}%
}
\newcommand{\defref}[2][]{%
	\hyperref[def:#2]{Definition~\ref{def:#2}#1}%
}

\newcommand{\Tfae}{The following are equivalent:}
\newcommand{\tfae}{the following are equivalent:}

% L-function on ξ — base field
\newcommand{\Lxi}{L\lt(s,\xi\rt)}
\newcommand{\ovLxi}{\ovL\lt(z,\xi\rt)}
% Lifted: \Lxil{k} = L^{(k)}(s,ξ^{(k)}), \ovLxil{k} = \ovL^{(k)}(z,ξ^{(k)})
\newcommand{\Lxil}[1]{L^{(#1)}\lt(s,\xi^{(#1)}\rt)}
\newcommand{\ovLxil}[1]{\ovL^{(#1)}\lt(z,\xi^{(#1)}\rt)}
% Custom first slot: \Lxis{arg} = L(arg,ξ), \Lxils{k}{arg} = L^{(k)}(arg,ξ^{(k)})
\newcommand{\Lxis}[1]{L\lt(#1,\xi\rt)}
\newcommand{\ovLxis}[1]{\ovL\lt(#1,\xi\rt)}
\newcommand{\Lxils}[2]{L^{(#1)}\lt(#2,\xi^{(#1)}\rt)}
\newcommand{\ovLxils}[2]{\ovL^{(#1)}\lt(#2,\xi^{(#1)}\rt)}

% Weil sums — base field
\newcommand{\Wpsi}[1]{W\lt(\psi;#1\rt)}
\newcommand{\Wchi}[1]{W\lt(\chi;#1\rt)}
\newcommand{\Wmix}[2]{W\lt(\psi,\chi;#1,#2\rt)}
% Lifted: \Wpsil{k}{f} = W^{(k)}(ψ^{(k)};f), etc.
\newcommand{\Wpsil}[2]{W^{(#1)}\lt(\psi^{(#1)};#2\rt)}
\newcommand{\Wchil}[2]{W^{(#1)}\lt(\chi^{(#1)};#2\rt)}
\newcommand{\Wmixl}[3]{W^{(#1)}\lt(\psi^{(#1)},\chi^{(#1)};#2,#3\rt)}
\newcommand{\sparsity}{\textit{sparsity}}

\let\oldchi\chi
\makeatletter
\DeclareRobustCommand{\chi}{\@ifnextchar_\sub@chi\latex@chi}
\newcommand{\latex@chi}{\protect\raisebox{2pt}{$\oldchi$}}
\newcommand{\sub@chi}[2]{%
  \@ifnextchar^{\subsup@chi{#2}}{\latex@chi^{}_{#2}}%
}
\newcommand{\subsup@chi}[3]{%
  \latex@chi_{#1}^{#3}%
}
\makeatother
\newcommand*\quot[2]{{^{\textstyle #1}\big/_{\textstyle #2}}}

\newenvironment{claimwidth}{\begin{center}\begin{adjustwidth}{0.05\textwidth}{0.05\textwidth}}{\end{adjustwidth}\end{center}}